% `\holdinginserts` (tex.web §1018). An insertion is moved into its box only
% when the parameter is zero or negative; positive means *hold*, and the
% insertion nodes stay in `\box255` exactly where they were.
%
% That is not a corner. It is how an output routine decides, *after* seeing the
% page, whether it wants the insertions on it: fire once with the parameter
% positive, look, `\unvbox255` to hand everything back — insertion nodes and
% all — set the parameter to zero, and let the builder consider them again. The
% second time they are placed for real. trip.tex's `\output` is built on
% exactly that, and LaTeX's float mechanism uses it to move a float to the next
% page after discovering it does not fit on this one.
%
% We had the parameter registered as a name and nothing read it, so insertions
% were always placed on the first firing. The second firing then saw none, and
% TeX's whole second page of trip.tex never happened here.
\catcode`\{=1 \catcode`\}=2 \catcode`\#=6
\font\testfont=cmr10 \testfont
\hsize=100pt \vsize=60pt \topskip=0pt \maxdepth=2pt \parindent=0pt
\dimen100=100pt \count100=500 \skip100=0pt

% Held: the box is still void when the routine looks, and the insertion is
% still in `\box255` to be handed back.
\holdinginserts=1
\output={\message{<fire hi=\the\holdinginserts box=\ifvoid100 void\else full\fi>}%
  \ifnum\holdinginserts>0
    \global\holdinginserts0 \unvbox255 \penalty-10000
  \else
    \message{<placed wd=\the\wd100 ht=\the\ht100>}%
    \shipout\box255
  \fi}
\hbox{A}\insert100{\hbox{X}}\penalty-10000

% Not held: the insertion is placed on the first firing, so the routine sees a
% full box straight away.
\global\holdinginserts=0
\output={\message{<direct box=\ifvoid100 void\else full\fi>}\shipout\box255}
\hbox{B}\insert100{\hbox{Y}}\penalty-10000

% Negative is the same as zero — §1018 tests `<= 0`.
\global\holdinginserts=-5
\output={\message{<negative box=\ifvoid100 void\else full\fi>}\shipout\box255}
\hbox{C}\insert100{\hbox{Z}}\penalty-10000

\end
